\documentclass[10pt]{beamer}

\usepackage{verbatim}
%\usepackage[latin1]{inputenc}
\usepackage{amssymb,amsmath,amsfonts,url}
\usepackage{graphicx,color}
\usepackage{pgf,pgfarrows,pgfnodes,pgfautomata,pgfheaps,pgfshade}
\usepackage[utf8]{inputenc}
\usepackage[french]{babel}
\usepackage{tikz}

\usepackage[all]{xy}
\usepackage{amstext}
\usepackage{xspace}
\usepackage{mdwtab}


\mode<article> {
  \usepackage{fullpage}
  \usepackage{hyperref}
  \setjobnamebeamerversion{advancedcourse.beamer}
}





\newcommand{\alertonce}[1]{\addtocounter{beamerpauses}{-1}\alert<+>{#1}}
\newcommand{\alerttwice}[1]{\addtocounter{beamerpauses}{-1}\alert<+-+(1)>{#1}}

%\newcommand{\subsubsubsection}[1]{\noindent\textbf{#1}}

\newcommand{\cfbox}[1]{\begin{center}\fbox{#1}\end{center}}
\newcommand{\cbox}[1]{\begin{center} #1 \end{center}}

%%%%%%%%%%%%%%%%%%%%%%%%%%%%%%%%%%%%%%%%%%%%%%%%%%%%%%%%%%%%%%%%%%%%%%%%%%%%%%%%%%%%%%%%%

\graphicspath{{Graphics/}}
%\input def_0.tex

%\input{StyleSlides.tex}
\mode<presentation> {
    \usepackage{beamerthemeBoadilla}

  %\beamertemplatetransparentcovereddynamic
  \beamertemplateballitem
}

\usepackage{listings}
\usepackage{xcolor}

\lstset{
	language=Python,
	basicstyle=\ttfamily\small,
	keywordstyle=\color{blue},
	commentstyle=\color{gray},
	stringstyle=\color{green!50!black},
	showstringspaces=false,
	frame=single,
	numbers=none
}
\lstset{
	language=Python,
	basicstyle=\ttfamily\small,
	keywordstyle=\color{blue},
	commentstyle=\color{gray},
	stringstyle=\color{green!50!black},
	showstringspaces=false,
	frame=single,
	numbers=none,
	literate=
	{é}{{\'e}}1
	{è}{{\`e}}1
	{ê}{{\^e}}1
	{ë}{{\"e}}1
	{à}{{\`a}}1
	{â}{{\^a}}1
	{ä}{{\"a}}1
	{ù}{{\`u}}1
	{û}{{\^u}}1
	{ü}{{\"u}}1
	{ô}{{\^o}}1
	{ö}{{\"o}}1
	{î}{{\^i}}1
	{ï}{{\"i}}1
	{ç}{{\c c}}1
}

\title[IN100 Cours 7]{\textbf{Fondements informatiques I}\\
Cours 7: Types et typage}

\author[]{Sorina Ionica \texttt{sorina.ionica@uvsq.fr} \\~\\Bilal FAYE  \texttt{bilal.faye@uvsq.fr} }


\institute{}

\date{}


%\setbeamercolor{normal text}{fg=black}

\begin{document}


\frame{\titlepage}
\begin{frame}[fragile]{Les séquences : définition et types}

Une \textbf{séquence} est un ensemble fini et ordonné
d'éléments. Chaque élément possède une position appelée
\textbf{indice}, qui commence à 0.

\medskip

\begin{center}
\begin{tikzpicture}[
    box/.style={
        draw,
        minimum width=0.8cm,
        minimum height=0.55cm
    }
]
\node[box] (a) at (0,0) {12};
\node[box] (b) at (0.8,0) {25};
\node[box] (c) at (1.6,0) {38};
\node[box] (d) at (2.4,0) {47};

\node at (0,-0.45) {\scriptsize 0};
\node at (0.8,-0.45) {\scriptsize 1};
\node at (1.6,-0.45) {\scriptsize 2};
\node at (2.4,-0.45) {\scriptsize 3};

\node at (1.2,-0.9) {\small Indices des éléments};
\end{tikzpicture}
\end{center}

En Python, on rencontre notamment quatre types de séquences :

\begin{itemize}
    \item \textbf{Listes} : modifiables (\texttt{list}).
    \item \textbf{Tuples} : non modifiables (\texttt{tuple}).
    \item \textbf{Chaînes de caractères} : non modifiables (\texttt{str}).
    \item \textbf{Intervalles} : séquences d'entiers (\texttt{range}).
\end{itemize}

\begin{exampleblock}{Exemples}
\begin{lstlisting}
l = [12, 25, 38, 47]
t = (12, 25, 38, 47)
s = "Python"
r = range(4)
\end{lstlisting}
\end{exampleblock}

\end{frame}
\begin{frame}[fragile]{Les séquences : opérations communes}

Les mêmes opérations peuvent être utilisées sur les
différents types de séquences.

\medskip

\small
\begin{center}
\begin{tabular}{|l|l|}
\hline
\textbf{Opération} & \textbf{Signification} \\
\hline
\texttt{len(s)} & Nombre d'éléments de \texttt{s} \\
\hline
\texttt{s[i]} & Élément d'indice \texttt{i} \\
\hline
\texttt{s[i:j]} & Éléments de \texttt{i} à \texttt{j-1} \\
\hline
\texttt{s[i:j:k]} & Même principe, avec un pas de \texttt{k} \\
\hline
\texttt{x in s} & Vérifie si \texttt{x} appartient à \texttt{s} \\
\hline
\texttt{s + t} & Concaténation de deux séquences \\
\hline
\texttt{s * n} & Répétition de la séquence \texttt{n} fois \\
\hline
\end{tabular}
\end{center}

\normalsize

\medskip

\begin{exampleblock}{Exemple en Python}
\begin{lstlisting}
s = [10, 20, 30, 40, 50]

print(s[0])       # 10
print(s[1:3])     # [20, 30]
print(len(s))     # 5
print(30 in s)    # True
\end{lstlisting}
\end{exampleblock}

\end{frame}
\begin{frame}[fragile]{Les séquences : parcours avec une boucle}

Les séquences sont des objets \textbf{itérables} :
on peut parcourir leurs éléments un par un.
La boucle \texttt{for} permet de réaliser ce parcours.

\begin{exampleblock}{Exemple en Python}
\begin{lstlisting}
# Une chaine de caracteres
ma_chaine = "Python"
for carac in ma_chaine:
    print(carac)
# Une liste
l = [55, 23, 103]
for nb in l:
    print(nb)
# Un tuple
t = (128, 0, 255)
for nb in t:
    print(nb)
\end{lstlisting}
\end{exampleblock}


\textbf{La boucle \texttt{for} parcourt les éléments
sans avoir besoin de manipuler leurs indices.}


\end{frame}
\begin{frame}[fragile]{Chaînes de caractères}

Les chaînes de caractères sont des séquences :
    on peut accéder à leurs caractères par indice et utiliser des tranches.


\begin{exampleblock}{Accès aux caractères et tranches}
\small
\begin{lstlisting}
chaine = 'Abracadabra!'
print(chaine[2])     # r
print(chaine[2:5])   # rac
print(chaine[-1])    # !
\end{lstlisting}
\end{exampleblock}

\textbf{Une particularité : les chaînes sont immuables}
\begin{itemize}
    \item Une chaîne ne peut pas être modifiée après sa création.
    \item On ne peut donc pas remplacer directement l'un de ses caractères.
\end{itemize}


\begin{exampleblock}{Une modification interdite}
\small
\begin{lstlisting}
chaine[2] = 'd'  # TypeError
liste = ["Paris", "Londres", "Rome"]
liste[1] = "Lille"       # Autorise : la liste est modifiable
liste[1][0] = 'l'        # TypeError : la chaine ne l'est pas
\end{lstlisting}
\end{exampleblock}

\end{frame}
\begin{frame}[fragile]{Des caractères aux nombres}

\begin{itemize}
    \item Un ordinateur manipule des nombres : pour représenter du texte,
    il faut associer un nombre à chaque caractère.
    \item L'\textbf{ASCII} est un premier système : il associe 128 caractères
    à des nombres de 0 à 127.
    \item \textbf{Unicode} étend ce principe à un très grand nombre de
    caractères (accents, alphabets, symboles, etc.).
\end{itemize}

\begin{block}{En Python : caractères et nombres}
    \small
    \begin{itemize}
        \item \texttt{ord(c)} donne le point de code Unicode du caractère \texttt{c}.
        \item \texttt{chr(n)} donne le caractère associé au point de code \texttt{n}.
    \end{itemize}
\end{block}

\begin{exampleblock}{Observer la correspondance}
\small
\begin{lstlisting}
print(ord('A'))  # 65
print(chr(65))   # 'A'
print(ord('é'))  # 233
print(chr(233))  # 'é'
\end{lstlisting}
\end{exampleblock}




\end{frame}
\begin{frame}[fragile]{Du texte aux octets}

\begin{itemize}
    \item En Python, une chaîne (\texttt{str}) est une séquence de caractères
    Unicode.
    \item Pour enregistrer ou transmettre du texte, il faut le convertir
    en octets : c'est le rôle de l'\textbf{encodage}.
\end{itemize}

\begin{exampleblock}{Encoder et décoder en UTF-8}
\small
\begin{lstlisting}
texte = "é"
octets = texte.encode("utf-8")
print(octets)  # b'\xc3\xa9'
texte_retrouve = octets.decode("utf-8")
print(texte_retrouve)  # é
\end{lstlisting}
\end{exampleblock}

\begin{center}
\small
\texttt{str} \quad $\xrightarrow{\texttt{encode}}$ \quad
\texttt{bytes} \qquad
\texttt{bytes} \quad $\xrightarrow{\texttt{decode}}$ \quad
\texttt{str}
\end{center}

\begin{itemize}
    \item UTF-8 est un encodage qui permet de représenter les caractères
    Unicode sous forme d'octets.
\end{itemize}

\end{frame}

\begin{frame}[fragile]{Les ensembles}

Un \textbf{ensemble (\texttt{set})} est une collection
d'éléments dans laquelle chaque élément apparaît une seule fois.

\begin{exampleblock}{Exemple}
\small
\begin{lstlisting}
A = {1, 3, 10, 5, 2}
B = {1, 10, 3, 3, 5, 2, 1}

print(A)
print(B)
print(A == B)  # True
\end{lstlisting}
\end{exampleblock}

\begin{itemize}
    \item Les doublons sont automatiquement supprimés.
    \item Les éléments ne sont pas ordonnés.
    \item Un ensemble est modifiable : on peut ajouter ou supprimer
    des éléments.
\end{itemize}

\end{frame}


\begin{frame}[fragile]{Créer des ensembles}

On peut créer un ensemble directement avec des accolades
ou à partir d'un objet itérable avec la fonction \texttt{set()}.

\begin{exampleblock}{Différentes façons de créer un ensemble}
\small
\begin{lstlisting}
# Avec des accolades
a = {1, 2, 3, 4}

# A partir d'une liste
b = set([1, 3, 4, 3, 1])

# A partir d'une chaine
c = set("Python")

# A partir de range()
d = set(range(5))
\end{lstlisting}
\end{exampleblock}

\begin{block}{Attention}
\small
Pour créer un ensemble vide, on utilise \texttt{set()}.

Les accolades vides \texttt{\{\}} créent un dictionnaire.
\end{block}

\end{frame}
\begin{frame}[fragile]{Parcourir un ensemble}

Un ensemble est \textbf{itérable} : on peut parcourir ses
éléments avec une boucle \texttt{for}.

En revanche, il n'est pas ordonné : on ne peut pas accéder
à un élément à partir de sa position.

\begin{exampleblock}{Parcourir un ensemble}
\small
\begin{lstlisting}
nombres = {10, 20, 30, 40}

for element in nombres:
    print(element)
\end{lstlisting}
\end{exampleblock}

\begin{alertblock}{Attention}
\small
L'instruction \texttt{nombres[0]} provoque une erreur :
un ensemble ne permet pas l'accès par indice.
\end{alertblock}

\end{frame}
\begin{frame}[fragile]{Supprimer les doublons}

Les ensembles sont particulièrement utiles pour supprimer
les doublons d'une liste.

\begin{block}{Problème}
\small
Écrire une fonction qui reçoit une liste et renvoie une
nouvelle liste contenant les mêmes éléments, sans doublons.
\end{block}

\pause

\begin{exampleblock}{Une solution avec les ensembles}
\small
\begin{lstlisting}
def supprime_doublons(liste):
    return list(set(liste))

nombres = [1, 3, 3, 5, 6, 7, 3, 1]

print(supprime_doublons(nombres))
\end{lstlisting}
\end{exampleblock}

\begin{alertblock}{Remarque}
\small
L'ordre des éléments n'est pas nécessairement conservé.
\end{alertblock}

\end{frame}
\begin{frame}{Opérations sur les ensembles}

Les ensembles permettent de réaliser plusieurs opérations
mathématiques.

Soient deux ensembles :
\[
A = \{0,1,2,3,4\}, \qquad
B = \{2,3,4,5\}
\]

\begin{itemize}
    \item \textbf{Union} : éléments appartenant à $A$ ou à $B$.
    \[
    A \cup B = \{0,1,2,3,4,5\}
    \]

    \item \textbf{Intersection} : éléments communs à $A$ et $B$.
    \[
    A \cap B = \{2,3,4\}
    \]

    \item \textbf{Différence} : éléments appartenant à $A$
    mais pas à $B$.
    \[
    A \setminus B = \{0,1\}
    \]

    \item \textbf{Différence symétrique} : éléments appartenant
    à un seul des deux ensembles.
    \[
    A \triangle B = \{0,1,5\}
    \]
\end{itemize}

\end{frame}
\begin{frame}[fragile]{Opérations sur les ensembles en Python}

Python propose des opérateurs pour réaliser les opérations
sur les ensembles.

\begin{center}
\begin{tabular}{c|c}
\textbf{Opérateur} & \textbf{Opération} \\
\hline
\texttt{|} & Union \\
\texttt{\&} & Intersection \\
\texttt{-} & Différence \\
\texttt{\^} & Différence symétrique
\end{tabular}
\end{center}

\begin{exampleblock}{Exemple}
\small
\begin{lstlisting}
A = {0, 1, 2, 3, 4}
B = {2, 3, 4, 5}

print(A | B)  # Union
print(A & B)  # Intersection
print(A - B)  # Difference
print(A ^ B)  # Difference symetrique
\end{lstlisting}
\end{exampleblock}

\end{frame}
\begin{frame}[fragile]{Les dictionnaires}

Un \textbf{dictionnaire (\texttt{dict})} est une structure de données
qui associe une \textbf{clé} à une \textbf{valeur}.

\begin{itemize}
    \item Chaque élément est un couple \texttt{clé : valeur}.
    \item On accède à une valeur grâce à sa clé, et non par sa position.
    \item Un dictionnaire est modifiable.
\end{itemize}

\begin{exampleblock}{Exemple : notes d'élèves}
\small
\begin{lstlisting}
notes = {
    "Gaetan": 14,
    "Laure": 9,
    "Kenza": 17,
    "Gabriel": 14
}

print(notes["Laure"])  # 9
\end{lstlisting}
\end{exampleblock}

\end{frame}
\begin{frame}[fragile]{Créer un dictionnaire}

On utilise des accolades \texttt{\{\}} pour créer un dictionnaire.
Chaque entrée est composée d'une clé et d'une valeur, séparées
par deux-points.

\begin{exampleblock}{Créer un dictionnaire}
\small
\begin{lstlisting}
eleve = {
    "nom": "Martin",
    "age": 18,
    "note": 15
}

print(eleve)
\end{lstlisting}
\end{exampleblock}

\begin{block}{Dictionnaire vide}
\small
On peut également créer un dictionnaire vide :
\begin{lstlisting}
dico = {}
\end{lstlisting}
\end{block}

\end{frame}
\begin{frame}[fragile]{Clés et valeurs}

Dans un dictionnaire :

\begin{itemize}
    \item Les \textbf{clés sont uniques} : une même clé ne peut
    désigner plusieurs entrées.
    \item Les \textbf{valeurs peuvent être répétées}.
    \item Les clés doivent être des objets \textbf{hachables},
    comme les entiers ou les chaînes de caractères.
\end{itemize}

\begin{exampleblock}{Exemple}
\small
\begin{lstlisting}
notes = {
    "Gaetan": 14,
    "Laure": 14,
    "Kenza": 17
}
\end{lstlisting}
\end{exampleblock}

\begin{alertblock}{Attention}
\small
Si une clé est utilisée plusieurs fois lors de la création
d'un dictionnaire, seule sa dernière valeur est conservée.
\end{alertblock}

\end{frame}
\begin{frame}[fragile]{Accéder aux valeurs}

Pour accéder à une valeur, on utilise sa clé entre crochets.

\begin{exampleblock}{Accéder à une valeur}
\small
\begin{lstlisting}
notes = {"Gaetan": 14, "Laure": 9, "Kenza": 17}

print(notes["Kenza"])  # 17
\end{lstlisting}
\end{exampleblock}

\begin{itemize}
    \item Si la clé existe, Python renvoie la valeur associée.
    \item Si la clé n'existe pas, Python déclenche une
    \texttt{KeyError}.
\end{itemize}

\begin{exampleblock}{Vérifier l'existence d'une clé}
\small
\begin{lstlisting}
if "Kenza" in notes:
    print(notes["Kenza"])
\end{lstlisting}
\end{exampleblock}

\end{frame}
\begin{frame}[fragile]{Modifier un dictionnaire}

Un dictionnaire est \textbf{modifiable} : on peut ajouter,
modifier ou supprimer des entrées.

\begin{exampleblock}{Ajouter et modifier une entrée}
\small
\begin{lstlisting}
notes = {"Gaetan": 14, "Laure": 9}

# Modifier une valeur existante
notes["Gaetan"] = 12

# Ajouter une nouvelle entrée
notes["Kenza"] = 17

print(notes)
\end{lstlisting}
\end{exampleblock}

\begin{block}{À retenir}
\small
\begin{itemize}
    \item Clé existante : la valeur est remplacée.
    \item Nouvelle clé : une entrée est ajoutée.
\end{itemize}
\end{block}

\end{frame}
\begin{frame}[fragile]{Parcourir un dictionnaire}

Une boucle \texttt{for} permet de parcourir les éléments
d'un dictionnaire.

\begin{exampleblock}{Parcourir les clés}
\small
\begin{lstlisting}
for cle in notes:
    print(cle)
\end{lstlisting}
\end{exampleblock}

\begin{exampleblock}{Parcourir les valeurs et les couples}
\small
\begin{lstlisting}
# Parcourir les valeurs
for valeur in notes.values():
    print(valeur)

# Parcourir les couples cle-valeur
for cle, valeur in notes.items():
    print(cle, ":", valeur)
\end{lstlisting}
\end{exampleblock}

\end{frame}
\begin{frame}[fragile]{Le typage dynamique}

En Python, le typage est \textbf{dynamique} :
le type d'une variable est déterminé automatiquement
à partir de la valeur qui lui est affectée.

\begin{exampleblock}{Exemple}
\small
\begin{lstlisting}
a = 5
print(type(a))  # <class 'int'>

a = "bonjour"
print(type(a))  # <class 'str'>
\end{lstlisting}
\end{exampleblock}

\begin{itemize}
    \item Une variable peut recevoir des valeurs de types différents.
    \item Il n'est pas nécessaire de déclarer son type à l'avance.
    \item Les opérations doivent néanmoins être compatibles avec
    les types des valeurs manipulées.
\end{itemize}

\begin{alertblock}{Attention}
\small
Le typage dynamique ne signifie pas que Python accepte
toutes les opérations entre types différents.
\end{alertblock}

\end{frame}

\begin{frame}[fragile]{Les opérations et les types}

Une opération doit être compatible avec les types
des valeurs manipulées.

\begin{exampleblock}{Exemples}
\small
\begin{lstlisting}
a = 5
b = 7
print(a + b)       # 12

x = "salut"
y = "hello"
print(x + y)       # saluthello

print(a + y)       # TypeError
print(y / 2)       # TypeError
\end{lstlisting}
\end{exampleblock}

\begin{itemize}
    \item Deux entiers peuvent être additionnés.
    \item Deux chaînes peuvent être concaténées.
    \item Certaines opérations entre types différents
    ne sont pas autorisées.
\end{itemize}

\end{frame}
\begin{frame}[fragile]{Les annotations de type}

Python permet d'ajouter des \textbf{annotations de type}
pour indiquer le type attendu d'une variable.

\begin{exampleblock}{Exemple}
\small
\begin{lstlisting}
age: int = 18
nom: str = "Alice"
taille: float = 1.72

notes: list[int] = [12, 15, 18]
\end{lstlisting}
\end{exampleblock}

\begin{itemize}
    \item \texttt{int} : entier.
    \item \texttt{float} : nombre décimal.
    \item \texttt{str} : chaîne de caractères.
    \item \texttt{list[int]} : liste d'entiers.
\end{itemize}

\begin{block}{À retenir}
\small
Les annotations sont facultatives : elles aident à comprendre
le programme, mais ne modifient pas le fonctionnement normal
de Python.
\end{block}

\end{frame}
\begin{frame}[fragile]{Annoter les paramètres d'une fonction}

On peut préciser le type attendu pour chaque paramètre
d'une fonction.

\begin{exampleblock}{Exemple}
\small
\begin{lstlisting}
def addition(a: int, b: int):
    return a + b

print(addition(5, 7))
\end{lstlisting}
\end{exampleblock}

\begin{itemize}
    \item \texttt{a: int} indique que le paramètre \texttt{a}
    est attendu comme un entier.
    \item \texttt{b: int} indique la même chose pour \texttt{b}.
\end{itemize}

\begin{alertblock}{Attention}
\small
Python ne vérifie pas automatiquement ces annotations
à l'appel de la fonction.

Elles servent notamment à guider le programmeur et les outils
de vérification de types.
\end{alertblock}

\end{frame}
\begin{frame}[fragile]{Annoter le type de retour}

On peut également indiquer le type de valeur renvoyé
par une fonction.

L'annotation \texttt{-> int} indique que la fonction
est censée renvoyer un entier.

\begin{exampleblock}{Exemple}
\small
\begin{lstlisting}
def somme(a: int, b: int) -> int:
    return a + b

resultat = somme(5, 7)
print(resultat)  # 12
\end{lstlisting}
\end{exampleblock}

\begin{block}{Vérification des types}
\small
Les annotations ne sont pas vérifiées automatiquement
par Python.

Un outil comme \texttt{mypy} permet d'analyser le programme
et de détecter certaines incohérences de types.
\end{block}

\end{frame}







\begin{comment}
\frame[containsverbatim]{
\frametitle{Les séquences}

\textcolor{blue}{\`A retenir} : Une séquence est ensemble fini et ordonné d'éléments indicés de $0$ à $n-1$ (si la séquence comporte $n$ éléments).

\begin{itemize}
\item les listes : modifiables
\item les tuples:  non-modifiables
\item les chaînes de caractères : non-modifiables
\item les "range" (intervalles)
\end{itemize}

\vspace{0.5 cm}

\small{
\begin{center}
\begin{tabular}{ |c|c| } 
 \hline
\texttt{x in s} &  \texttt{True} si un élément de \texttt{s} est égal à \texttt{x}, sinon \texttt{False}  \\ 
\texttt{x not in s} & \texttt{False} si un élément de \texttt{s} est égal à \texttt{x}, sinon \texttt{True} \\ 
\texttt{s + t}  &  la concaténation de \texttt{s} et \texttt{t}  \\ 
\texttt{s * n} ou \texttt{n * s} & ajouter \texttt{s} $n$ fois à lui-même \\
\texttt{s[i]} & $i$-ème élément de \texttt{s} en commençant par 0 \\
\texttt{s[i:j]} & tranche (slice) de \texttt{s} de \texttt{i} à \texttt{j-1} \\
\texttt{s[i:j:k]} & tranche (slice) de \texttt{s} de \texttt{i} à \texttt{j-1} avec un pas de \texttt{k} \\
\texttt{len(s)} &  longueur de \texttt{s} \\
 \hline
\end{tabular}
\end{center}}


}

\frame[containsverbatim]{
\frametitle{Les séquences}

\textcolor{blue}{\`A retenir} : Les séquences sont des objets \textbf{itérables.}\\


 On peut  boucler sur tous éléments grâce à une boucle \textbf{for}.



\begin{exampleblock}{}
\small{
\begin{verbatim}
ma_chaine = "Abracadabra"

for carac in ma_chaine:
  print(carac)
  
l = [55, 23,103,258]
t = (128, 0, 255)

for nb in l:
    print(nb)

for nb in t:
    print(nb)
 \end{verbatim}
 }
\end{exampleblock}


}
\begin{frame}[fragile]{Les séquences : parcours avec une boucle}

Les séquences sont des objets \textbf{itérables} :
on peut parcourir leurs éléments un par un.

\medskip

La boucle \texttt{for} permet de réaliser ce parcours.

\begin{exampleblock}{Exemple en Python}
\begin{lstlisting}
# Une chaine de caracteres
ma_chaine = "Python"

for carac in ma_chaine:
    print(carac)

# Une liste
l = [55, 23, 103]

for nb in l:
    print(nb)

# Un tuple
t = (128, 0, 255)

for nb in t:
    print(nb)
\end{lstlisting}
\end{exampleblock}

\medskip

\begin{center}
\textbf{La boucle \texttt{for} parcourt les éléments
sans avoir besoin de manipuler leurs indices.}
\end{center}

\end{frame}
\begin{frame}[fragile]{Les séquences : parcours avec une boucle}

Les séquences sont des objets \textbf{itérables} :
on peut parcourir leurs éléments un par un.


\end{frame}


\begin{frame}[fragile]{Chaînes de caractères}


\begin{itemize}
%\item  Les chaînes de caractères peuvent être vues comme une sorte de listes particulières.
\item 
On peut notamment utiliser  les \emph{tranches}.
\end{itemize}

\begin{exampleblock}{}
\small{
\begin{verbatim}
chaine = 'Abracadabra!'
print(chaine[2:5])
print(chaine[-1])
\end{verbatim}
}
\end{exampleblock}

\vspace{1 cm}

\begin{itemize}
\item  La différence principale entre les listes et les chaînes de caractères
est que les chaînes de caractères sont des objets
\textbf{non-modifiables} après création.
\end{itemize}

\begin{exampleblock}{}
\small{
\begin{verbatim}
chaine[2] = 'd'  #TypeError: 'str' object does not support item assignment 
liste_chaines=["Paris","londres", "Rome","Budapest"]
liste_chaines[1][0]='L'
\end{verbatim}
}
\end{exampleblock}

\end{frame}

\begin{frame}[fragile]{Les encodages}

\small{En Python, un caractère est une chaîne de longueur 1. \\}

\begin{block}{\small{L'encodage \textbf{ASCII} (American Standard Code for Information Interchange)}}
\small{
 ASCII représente 128 caractères anglais sous forme de nombres, chaque lettre étant associée à un numéro spécifique compris entre 0 et 127. 
 \begin{itemize}
 \item Par exemple, le code ASCII pour la lettre majuscule A est 65, pour la majuscule B c’est 66 etc. 
 \end{itemize}
 }
\end{block}    
    

\small{
\begin{itemize}
\item La fonction \texttt{ord} convertit un caractère dans sa représentation numérique. 
\item La fonction \texttt{chr} convertit un nombre dans un caractère ASCII. 
\end{itemize}
}

\pause
\vspace{0.5 cm}
\small{
En fait, Python repose sur l'encodage Unicode, qui comprend beaucoup plus de caractères: 
\begin{exampleblock}{}
\begin{verbatim}
>>> chr(245)
'õ'
\end{verbatim}
\end{exampleblock}
}
\end{frame}

\begin{frame}[fragile]{Les ensembles}

Un \textbf{ensemble (set)} est une collection d'éléments ne contenant
pas de doublons.

\vspace{0.3 cm}

\textbf{Exemple:}  $A = \{1, 3, 10, 5, 2\}$ est un ensemble, tandis que
$B = \{1, 10, 3, 3, 5, 2, 1\}$ n'en est pas un.

\vspace{0.3 cm}

Python permet de créer des objets de type \texttt{set}.

\begin{itemize}
\item Un \texttt{set} n'est ni \textbf{ordonné}, ni \textbf{modifiable}.
\item On utilise des \texttt{\{\}} pour créer un nouveau \texttt{set}.
\end{itemize}

\vspace{0.3 cm}
\begin{exampleblock}{}
\small{
\begin{verbatim}
s = {1, 2, 3, 4, 5}
print(s)
print(type(s))
\end{verbatim}
}
\end{exampleblock}
\end{frame}




\frame[containsverbatim]{
\frametitle{Créer des ensembles}


La fonction \texttt{set()} permet de générer un ensemble à partir de
n'importe quel objet itérable.

\vspace{0.5 cm}

\begin{exampleblock}{}
\begin{verbatim}
# À partir d'une liste
s = set([1, 3, 4, 5, 4, 3])

# À partir de range()
s = set(range(10))

# À partir d'une chaîne de caractères
s = set("Versailles")
\end{verbatim}
\end{exampleblock}

}


\frame[containsverbatim]{
\frametitle{Accéder aux éléments d'un ensemble}

\begin{itemize}
\item  Un \texttt{set} est \textbf{itérable}. 
 \item Un set n'est pas \textbf{ordonné}, donc il n'est pas une séquence. 
 \begin{itemize}
\item  Il est impossible de récupérer un élément par sa position.
\end{itemize}
\end{itemize}

\begin{exampleblock}{}
\begin{verbatim}
s = {0, 5, 10, 15}
for element in s :
    print(element)
    
print(s[2])     #TypeError: 'set' object is not subscriptable
\end{verbatim}
\end{exampleblock}


}
  
  
\begin{frame}[fragile]{Doublons}  


Les \texttt{sets} sont très pratiques pour \emph{supprimer des
doublons}.

\vspace{0.5 cm}

\textbf{Exercice} : Écrire une fonction qui prend en paramètre une liste
et renvoie une nouvelle liste, égale à celle passée en paramètre mais
sans les doublons.

\pause

\vspace{0.5 cm}

\begin{exampleblock}{}
\begin{verbatim}
def supprime_doublons(liste) :
    return list(set(liste))

nouvelle_liste = supprime_doublons([1, 3, 3, 5, 6, 7, 3, 1])
print(nouvelle_liste)
\end{verbatim}
\end{exampleblock}


\end{frame}




\frame[containsverbatim]{
\frametitle{Opérateurs sur les ensembles}


Les ensembles permettent d'effectuer des opérations mathématiques telles
que l'\emph{union}, l'\emph{intersection}, la \emph{différence} ou la
\emph{différence symétrique}.

\vspace{0.5 cm}

\textbf{Exemple:} \(A = \{0, 1, 2, 3, 4\}\), \(B = \{2, 3, 4, 5\}\)

\begin{itemize}
\item
  \(A\cup B := \{x \in A \text{ ou } x\in B\} = \{0, 1, 2, 3, 4, 5\}\)
\item
  \(A\cap B := \{x \in A \text{ et } x\in B\} = \{2, 3, 4\}\)
\item
  \(A\setminus B := \{x \in A \text{ et } x \not \in B\} = \{0, 1\}\)
\item
  \(A\triangle B := \{x \in A\cup B \text{ et } x \not \in A\cap B\} = \{0, 1, 5\}\)
\end{itemize}

}


\frame[containsverbatim]{
\frametitle{Opérateurs sur les ensembles}



 \begin{itemize}
\item
  \texttt{a\ \textbar{}\ b}: \(A\cup B\)
\item
  \texttt{a\ \&\ b}: \(A\cap B\)
\item
  \texttt{a\ -\ b}: \(A \setminus B\)
\item
  \texttt{a\ \^{}\ b}: \(A\triangle B\)
\end{itemize}


\begin{exampleblock}{}
\begin{verbatim}
chaine1 = 'chaussette'
chaine2 = 'archiduchesse'

a = set(chaine1)
b = set(chaine2)

print(a | b)
print(a & b)
print(a - b)
print(a ^ b)
\end{verbatim}
\end{exampleblock}




}
\frame[containsverbatim]{
\frametitle{Dictionnaires}


Un dictionnaire est une structure de données, dont les éléments sont
accessibles grâce à une \emph{clé}.

\vspace{0.5 cm}

 \begin{itemize}
\item Pour construire un dictionnaire, on utilise des accolades \texttt{\{\}}.
\item Un élément d'un dictionnaire est un couple \texttt{cle:\ valeur}.
\item On accède à une \textbf{valeur} du dictionnaire en utilisant la \textbf{clé}. 
\end{itemize}


\vspace{0.5 cm}
\begin{exampleblock}{Exemple}
\small{
\begin{verbatim}
dico_notes = {"Gaëtan" : 14, "Laure" : 9, "Kenza" : 17, "Gabriel" : 14}

print(dico_notes["Laure"])
\end{verbatim}
}
\end{exampleblock}

}


\frame[containsverbatim]{
\frametitle{Dictionnaires}



\begin{itemize}
\item Dans un dictionnaire, les clés doivent être \emph{uniques}, mais pas
  les valeurs.
\item Les clés peuvent être n'importe quel objet non-modifiable (immutable) 
\begin{itemize}
\item Par example des \texttt{int} ou des \texttt{string}. 
\end{itemize}
\item Il est facile de créer et modifier un dictionnaire.
\end{itemize}



\vspace{0.5 cm}

\begin{exampleblock}{}
\small{
\begin{verbatim}
dico_notes = {"Gaëtan" : 14, "Laure" : 9, "Kenza" : 17, "Gabriel" : 14}

# Changer une entrée existante
dico_notes["Gaëtan"] = 12

# Ajouter une nouvelle entrée

dico_notes["Anne"] = 15
print(dico_notes)

# Créer un dictionnaire vide.
dico_vide = {}
\end{verbatim}
}
\end{exampleblock}


}



%# Utilisation du constructeur dict()
%dico_restos = dict([("L'auberge du Goëllo", 22170), ("Café Menilmontant", 75011), ("Mamie Bigoude", 17000)])
%print(dico_restos)


\frame[containsverbatim]{
\frametitle{Parcourir un dictionnaire}

On peut vérifier facilement qu'une valeur est dans le dictionnaire si on connaît sa clé. 


\begin{exampleblock}{}
\begin{verbatim}
cle="Kenza"
if cle in dico_notes:
     print(cle, "a obtenu la note", dico_notes[cle])     
\end{verbatim}
\end{exampleblock}


\vspace{0.5 cm}

Les clés et leurs valeurs peuvent être récupérées en même temps en
  utilisant la méthode \texttt{items()}. 

\vspace{0.5 cm}

\begin{exampleblock}{}
\small{
\begin{verbatim}
dico_notes = {"Gaëtan" : 14, "Laure" : 9, "Kenza" : 17, "Gabriel" : 14}

for i in dico_notes.items():
    print(i)
\end{verbatim}
}
\end{exampleblock}


}

\begin{frame}[fragile]{Typage en Python}


\textbf{Rappel} : En Python le typage est dynamique. 

\begin{exampleblock}{}
\begin{verbatim}
a=5
b=7
print(a+b)  # int + int
x="salut"
y="hello"
print("salut" + "hello") # str + str
print(a + y)            #int + str   Erreur!

# Problème de type pour l'opération /
"hello" / 2
\end{verbatim}
\end{exampleblock}
\end{frame}


\begin{frame}[fragile]{Annotation de type}



\begin{itemize}
\item Depuis la version 3.5 on peut ajouter \textbf{des annotations de type} aux variables et fonctions. 
\item Ces annotations  permettent d'apporter plus de structure aux programmes pour éviter certains bugs. 
\item Ces annotations sont facultatives et vous pouvez mélanger du code annoté ou non. 
\end{itemize}

\vspace{0.5 cm}


\begin{exampleblock}{}
\begin{verbatim}
a: int = 5
b: int =7
print(a+b)  # int + int
x: str="salut"
y: str = "hello"

y=6
print(type(y))

liste: list[int] = [1, 2, 3]
\end{verbatim}
\end{exampleblock}

\end{frame}


\begin{frame}[fragile]{Annotation de type pour les fonctions}



On peut spécifier le type des arguments. 
 \begin{itemize}
 \item  Les annotations \texttt{a: int}, \texttt{b:int} indiquent que les paramètres doivent être des entiers. 
 \end{itemize}

\begin{exampleblock}{Exemple}
\small{
\begin{verbatim}
def mon_print(a : int, b: int):     
	       print(a + b)
      
mon_print(5,7)
mon_print("salut","hello")
mon_print(5,"salut")
\end{verbatim}
}
\end{exampleblock}


\begin{itemize}
\item Les informations de type ne changent rien à l'exécution. 
\item  Il y a une erreur de type seulement  quand on applique une fonction à un élément et qu'il n'est pas du type attendu par la fonction.
\end{itemize}

\end{frame}


\begin{frame}[fragile]{Annotation de type}



On peut spécifier le type de retour. 
\begin{itemize}   
\item    L'annotation  \texttt{-> int} indique que la fonction renvoie un entier.
 \end{itemize}


\begin{exampleblock}{Exemple}
\small{
\begin{verbatim}
def somme(a: int, b: int) -> int:
	      return a + b


print(somme(5,7))
print(somme("salut","hello"))	
\end{verbatim}
}
\end{exampleblock}

\begin{itemize}
\item Pour détecter les erreurs, il faut utiliser des outils comme \texttt{mypy}. 
\end{itemize}



\end{frame}
\end{comment}
\end{document}

